% Copyright (c) 2026 Geovana Neves. Licensed under CC BY 4.0 (https://creativecommons.org/licenses/by/4.0/). See LICENSE-AND-AUTHORSHIP.md.
%
% 07-python-environment-offline/main.tex
%
% Guide 7 of the pyflightstream guides: Python, a virtual environment and
% pyflightstream on a machine with no internet, on Windows and on Linux.
% Built by ../build-all.ps1 (or build-all.sh) with this folder and ../shared/
% as include directories.
%
% The download commands were run on 2026-09-28 against the package index:
% for Linux, manylinux_2_28_x86_64 ALONE fails (pydantic-core publishes only
% manylinux2014 wheels) and manylinux2014_x86_64 ALONE fails with the fsi
% extra (PyNiteFEA needs numpy >= 2.4, which publishes no manylinux2014
% wheel); the two tags together resolve.
%
% The sentence "To run a matrix locally on Linux, change only the executable
% path in executables.toml." is kept verbatim, whole and unhyphenated.
%
% No em dashes and no en dashes. No underscores in frame or section titles
% outside \code: beamer writes titles to .nav and .toc.

\documentclass[aspectratio=169,11pt]{beamer}

\input{preamble-slides.tex}
\input{hooks-slides.tex}
\ftsinput{boxes}
\ftsinput{hooks-contract}

\usepackage[nohyphen]{underscore}

\input{info.tex}

\begin{document}

\guidetitleframe{7}{Python environment for offline machines}{Offline Python environment}
  {Python, a virtual environment and \pkgname{} on a machine with no internet, Windows and Linux}
  {Sources: the package's installation notes and the \code{pip} documentation [1, 2].}

\begin{frame}{Who reads this}
\begin{itemize}
  \item The person who installs \pkgname{} on a machine that cannot reach the
        internet: a Windows workstation behind a closed network, a Linux
        workstation, or the shared file system of an HPC cluster.
  \item It happens once per machine and per version. Everybody else activates
        the environment with one line (Guide 1).
  \item The same recipe serves both systems: \textbf{download on a connected
        machine, carry one folder, install with the network forbidden}.
\end{itemize}
\end{frame}

\begin{frame}{Contents}
\small
\begin{multicols}{2}
\tableofcontents
\end{multicols}
\end{frame}

\section{The idea}

\begin{frame}{Two machines and one folder}
\begin{center}
\begin{tikzpicture}[font=\scriptsize, node distance=9mm,
  m/.style={draw=CourseBlue, thick, rounded corners, fill=CourseBlue!8,
            align=center, text width=36mm, minimum height=16mm, inner sep=3pt},
  f/.style={draw=CourseGold, thick, rounded corners, fill=CourseGold!15,
            align=center, text width=28mm, minimum height=12mm, inner sep=3pt}]
\node[m] (on) {\textbf{Connected machine}\\ any system, any Python\\ \tiny \code{pip download} for the TARGET};
\node[f, right=of on] (w) {\textbf{one folder of wheels}\\ \tiny every file ends in \code{.whl}};
\node[m, right=of w] (off) {\textbf{Offline machine}\\ Windows or Linux\\ \tiny \code{pip install --no-index}};
\draw[-{Stealth}, thick, CourseBlue] (on) -- (w);
\draw[-{Stealth}, thick, CourseBlue] (w) -- node[above, font=\tiny] {carry} (off);
\end{tikzpicture}
\end{center}
\begin{itemize}\small
  \item A \textbf{wheel} is a ready-built package. With wheels only, the
        offline machine needs no compiler and no network.
  \item The download names the \textbf{target's} Python version and platform,
        not the connected machine's.
  \item The install forbids the network, so it either installs from what was
        carried or fails naming what is missing. It never hangs.
\end{itemize}
\end{frame}

\begin{frame}{What the target needs to be}
\begin{columns}[T]
\begin{column}{0.48\textwidth}
\begin{block}{Python 3.12 or newer}
\begin{itemize}\small
  \item The package declares \code{requires-python = ">=3.12"}, following the
        scientific-python support schedule SPEC 0 [3].
  \item 3.9 and 3.10 lack \code{tomllib} and \code{enum.StrEnum}, which the
        package uses. 3.11 is below the declared floor: \code{pip} refuses it
        naming \code{Requires-Python}.
  \item Download for the same minor version the target runs: \code{312}.
\end{itemize}
\end{block}
\end{column}
\begin{column}{0.48\textwidth}
{\ftstablesize
\begin{tabular}{@{}ll@{}}
\toprule
\textbf{Target} & \textbf{Platform tags to download} \\
\midrule
Windows, 64 bit & \code{win\_amd64} \\
Linux x86-64,   & \code{manylinux\_2\_28\_x86\_64} \\
glibc 2.28 or later & \emph{and} \code{manylinux2014\_x86\_64} \\
\bottomrule
\end{tabular}}
\vspace{2mm}
\begin{alertblock}{Linux: check the C library first}
\small \code{ldd --version} on a node must print 2.28 or later (RHEL 8
and its rebuilds, and newer). Older systems cannot run the current
\code{numpy} wheels [4].
\end{alertblock}
\end{column}
\end{columns}
\end{frame}

\begin{frame}{Why two Linux tags}
\begin{itemize}\small
  \item \code{--platform} tells \code{pip} which wheels the target accepts. It
        accepts \textbf{exactly the tags you list}, so list every one the
        target can run.
  \item The structural solver of the FSI extra, PyNiteFEA, needs
        \code{numpy >= 2.4}, whose Linux wheels are \code{manylinux\_2\_28}
        only: with \code{manylinux2014} alone the download fails.
  \item Other dependencies (\code{pydantic-core}) publish \code{manylinux2014}
        wheels only: with \code{manylinux\_2\_28} alone the download fails
        too.
  \item Both tags together resolve. A glibc 2.28 system runs both kinds.
\end{itemize}
\begin{takeaway}If the download fails with \emph{ResolutionImpossible}, a
platform tag is missing; do not drop \code{--only-binary}.\end{takeaway}
\end{frame}

\section{Python on the offline machine}

\begin{frame}[fragile]{Windows: one installer, no administrator}
\begin{itemize}\small
  \item On the connected machine, download the \textbf{Windows installer
        (64-bit)} of the latest Python 3.12 from \code{python.org} [5], and
        carry it with the wheels.
  \item On the offline machine, run it for the current user only: leave
        \emph{Install for all users} off. No administrator is needed.
  \item Keep the \code{py} launcher on. Check:
\end{itemize}
\begin{lstlisting}[style=ftsbase]
py -3.12 --version
\end{lstlisting}
\begin{alertblock}{Corporate images}
\small A machine may already carry a managed Python. Use it only if it is 3.12
or newer; otherwise install beside it, per user, and always call it as
\code{py -3.12}.
\end{alertblock}
\end{frame}

\begin{frame}{Linux: three routes to a 3.12, no administrator}
\slidefig[0.46]{guide-three-routes}
\begin{itemize}\small
  \item A cluster's system Python is usually older than 3.12. Put a 3.12
        \textbf{inside the shared install folder} (\code{\clusterroot}); none
        of the routes touches the system Python.
  \item Route C always works on an offline machine: a relocatable CPython from
        the python-build-standalone project [6], copied in.
\end{itemize}
\end{frame}

\begin{frame}[fragile]{Linux, Route C: choosing the standalone archive}
% \relax first: a brace group right after the title is read as a subtitle.
\relax
{\ftstablesize
\begin{tabular}{@{}lll@{}}
\toprule
\textbf{Choice} & \textbf{Take} & \textbf{Why not the other} \\
\midrule
version   & \code{3.12}       & the wheels are \code{cp312}; a 3.13 rejects them \\
CPU level & \code{x86\_64}    & \code{x86\_64\_v3} needs AVX2 and dies on an older node \\
contents  & \code{install\_only\_stripped} & the same interpreter, about a third of the size \\
libc      & \code{linux-gnu}  & \code{musl} is for Alpine \\
\bottomrule
\end{tabular}}\par
\vspace{2mm}
\small Copy the archive into the shared folder, then on the Linux machine:
\begin{lstlisting}[style=ftsbase]
export FTS=(*@\clusterroot@*)
cd $FTS
sha256sum cpython-3.12*-install_only_stripped.tar.gz   # compare with the release
tar -xzf  cpython-3.12*-install_only_stripped.tar.gz   # creates $FTS/python
$FTS/python/bin/python3.12 --version
\end{lstlisting}
\small The interpreter finds its standard library relative to its own binary
and writes nothing outside \code{\$FTS}. Nothing is compiled.
\end{frame}

\section{The wheels}

\begin{frame}[fragile]{Download on the connected machine}
\small For a Linux target (one command; \code{\^{}} continues a line in the
Windows shell, \code{\textbackslash} in a Linux one):
\begin{lstlisting}[style=ftsbase]
python -m pip download "pyflightstream[fsi]==0.30.0" -d (*@\envname@*)-linux ^
    --python-version 312 --only-binary=:all: ^
    --platform manylinux_2_28_x86_64 --platform manylinux2014_x86_64
\end{lstlisting}
\small For a Windows target:
\begin{lstlisting}[style=ftsbase]
python -m pip download "pyflightstream[fsi]==0.30.0" -d (*@\envname@*)-win ^
    --python-version 312 --only-binary=:all: --platform win_amd64
\end{lstlisting}
{\ftstablesize
\begin{tabular}{@{}ll@{}}
\toprule
\textbf{Option} & \textbf{What it prevents} \\
\midrule
\code{--python-version 312} & downloading for the connected machine's Python \\
\code{--platform ...}       & downloading the connected machine's own wheels \\
\code{--only-binary=:all:}  & downloading source, which would need a compiler \\
\bottomrule
\end{tabular}}
\end{frame}

\begin{frame}[fragile]{A development build, and the check before carrying}
\begin{columns}[T]
\begin{column}{0.50\textwidth}
\small A development build (a version ending in \code{.dev<N>}) is not on
the package index. Name its wheel file instead of the version; \code{pip}
copies it and downloads its dependencies:
\begin{lstlisting}[style=ftsbase,basicstyle=\ttfamily\tiny]
python -m pip download ^
  "pyflightstream-<version>.dev<N>-py3-none-any.whl[fsi]" ^
  -d (*@\envname@*)-win --python-version 312 ^
  --only-binary=:all: --platform win_amd64
\end{lstlisting}
\end{column}
\begin{column}{0.46\textwidth}
\begin{alertblock}{The failure that gives no warning}
\small Without \code{--platform} and \code{--python-version}, \code{pip}
quietly downloads for \emph{your} machine. The folder looks right and copies
fine, and the install fails on the target after everything was moved.
\end{alertblock}
\begin{block}{Check}
\small Every file must end in \code{.whl}; one \code{.tar.gz} means source got
in. The platform files must carry the target's tag.
\end{block}
\end{column}
\end{columns}
\end{frame}

\section{Installing}

\begin{frame}[fragile]{Linux: the environment at its final path}
\begin{lstlisting}[style=ftsbase]
$FTS/python/bin/python3.12 -m venv --copies $FTS/py-venv/(*@\envname@*)
source $FTS/py-venv/(*@\envname@*)/bin/activate
ls $FTS/pfs-wheels/(*@\envname@*) | wc -l      # the folder pip will read
pip install --no-index --find-links $FTS/pfs-wheels/(*@\envname@*) "pyflightstream[fsi]"
python -c "import pyflightstream; print(pyflightstream.__version__)"
pyfs-matrix --help
\end{lstlisting}
\begin{itemize}\small
  \item \code{--no-index} forbids the internet; \code{--find-links} says where
        to look instead.
  \item \code{--copies} puts the interpreter inside the environment, so it
        survives anybody tidying up \code{python/} later.
  \item The version line must print the version you downloaded, and
        \code{pyfs-matrix --help} must answer.
\end{itemize}
\end{frame}

\begin{frame}[fragile]{Windows: the same, with the launcher}
\begin{lstlisting}[style=ftsbase]
cd (*@\mainworkspace@*)
py -3.12 -m venv .venv
.venv\Scripts\activate
pip install --no-index --find-links <folder>\(*@\envname@*)-win "pyflightstream[fsi]"
python -c "import pyflightstream; print(pyflightstream.__version__)"
pyfs-matrix --help
\end{lstlisting}
\begin{itemize}\small
  \item \code{<folder>} is where the carried wheels are; any local or network
        folder works.
  \item \code{deactivate} leaves the environment.
  \item A development build installs the same way: its wheel is in the folder.
\end{itemize}
\end{frame}

\begin{frame}{Rules of a shared environment}
\begin{columns}[T]
\begin{column}{0.48\textwidth}
\begin{alertblock}{Build it at its final path}
\small A virtual environment cannot be moved: its absolute path is written
into \code{activate} and into every command. Built in \code{\$HOME} and copied,
it breaks for a colleague days later with \code{bad interpreter}.
\end{alertblock}
\end{column}
\begin{column}{0.48\textwidth}
\begin{block}{The interpreter is recorded, not configured}
\small A run writes the path of the Python that ran it into the solver's
action programs. The environment must therefore sit at a path valid on every
node.
\end{block}
\end{column}
\end{columns}
\vspace{2mm}
\begin{takeaway}Put the version in the folder name (\code{\envname}). Older
environments stay beside the new one, and an old study stays
reproducible.\end{takeaway}
\end{frame}

\begin{frame}[fragile]{Two folders with the same name}
\small \code{py-venv/\envname} is the environment, and \code{venv} creates it.
\code{pfs-wheels/\envname} is the folder of wheels, and only the copy from the
connected machine creates it. When the second is missing, \code{pip} answers:
\begin{lstlisting}[style=ftsbase]
WARNING: Location '.../pfs-wheels/(*@\envname@*)' is ignored: it is either
a non-existing path or lacks a specific scheme.
ERROR: No matching distribution found for pyflightstream
\end{lstlisting}
\begin{takeaway}An \code{ls} of \code{py-venv/} shows a folder that looks like
proof. It is the other folder: list the one \code{pip} reads.\end{takeaway}
\end{frame}

\begin{frame}[fragile]{One activation file for everybody (Linux)}
\begin{lstlisting}[style=ftsbase]
# $FTS/bin/(*@\envname@*)-env
export PATH=(*@\clusterroot@*)/py-venv/(*@\envname@*)/bin:$PATH
\end{lstlisting}
\begin{itemize}\small
  \item Users write \code{source \$FTS/bin/\envname-env} at the top of a
        session or of a job script.
  \item It exports \code{PATH} rather than defining an alias, because a job
        script does not expand aliases.
  \item The files of older environments stay beside it.
  \item To leave it: a new shell, or take the folder off \code{PATH} again.
\end{itemize}
\end{frame}

\section{Running and recording}

\begin{frame}[fragile]{Running a matrix locally on Linux}
\begin{block}{The rule}
{\hyphenpenalty=10000\exhyphenpenalty=10000\relax
\textbf{To run a matrix locally on Linux, change only the executable path in executables.toml.}\par}
\end{block}
\begin{itemize}\small
  \item The row's \code{FS\_BUILD} id maps to the executable in
        \code{inputs/executables.toml}; a path true of one machine only goes in
        \code{inputs/executables.local.toml}, read over it. Nothing else changes.
  \item With a profile in \code{inputs/hpc/} the run would submit; keep it here:
\end{itemize}
\begin{lstlisting}[style=ftsbase]
source $FTS/bin/(*@\envname@*)-env
pyfs-matrix run matriz.fs --workspace . --local
\end{lstlisting}
\small Records say \code{forced\_local}; no \code{collect} needed afterwards.
\end{frame}

\begin{frame}{Record what you did}
\relax
{\ftstablesize
\begin{tabular}{@{}>{\raggedright\arraybackslash}p{0.36\linewidth}>{\raggedright\arraybackslash}p{0.56\linewidth}@{}}
\toprule
\textbf{Fact} & \textbf{Why the next person needs it} \\
\midrule
where the 3.12 came from & the module name, or the standalone archive and its sha256 \\
the environment folder   & \code{py-venv/\envname}, and the activation file \\
the package version      & what \code{pyflightstream.\_\_version\_\_} printed \\
the platform tags used   & the download command, verbatim \\
when, and by whom        & so the next installation is the same one \\
\bottomrule
\end{tabular}}
\vspace{2mm}
\begin{takeaway}Two facts exist only in the head of whoever installed. Write
them down, beside the environment.\end{takeaway}
\end{frame}

\begin{frame}[fragile]{Command summary}
\begin{columns}[T]
\begin{column}{0.50\textwidth}
\small\textbf{Windows target}
\begin{lstlisting}[style=ftsbase,basicstyle=\ttfamily\tiny]
pip download "pyflightstream[fsi]==0.30.0" ^
  -d W --python-version 312 ^
  --only-binary=:all: --platform win_amd64
py -3.12 -m venv .venv
.venv\Scripts\activate
pip install --no-index --find-links W ^
  "pyflightstream[fsi]"
pyfs-matrix --help
\end{lstlisting}
\end{column}
\begin{column}{0.46\textwidth}
\small\textbf{Linux target}
\begin{lstlisting}[style=ftsbase,basicstyle=\ttfamily\tiny]
pip download "pyflightstream[fsi]==0.30.0" \
  -d L --python-version 312 \
  --only-binary=:all: \
  --platform manylinux_2_28_x86_64 \
  --platform manylinux2014_x86_64
python3.12 -m venv --copies $FTS/py-venv/E
source $FTS/py-venv/E/bin/activate
pip install --no-index --find-links L \
  "pyflightstream[fsi]"
\end{lstlisting}
\end{column}
\end{columns}
\end{frame}

\begin{frame}{References}
\begin{reflist}
  \bibitem{g07r01} \docref{getting-started}{Getting started}, and the repository's
        \code{README.md}: \url{\pkgrepo}.
  \bibitem{g07r02} Python Packaging Authority, \emph{pip documentation}: \code{pip
        download} and \code{pip install}, 2026.
        \url{https://pip.pypa.io/en/stable/cli/pip_download/}
  \bibitem{g07r03} Scientific Python, \emph{SPEC 0: minimum supported dependencies}.
        \url{https://scientific-python.org/specs/spec-0000/}
  \bibitem{g07r04} N. J. Smith, T. Kluyver, \emph{PEP 600: future manylinux platform
        tags for portable Linux built distributions}, 2019.
        \url{https://peps.python.org/pep-0600/}
  \bibitem{g07r05} Python Software Foundation, \emph{Python releases for Windows}.
        \url{https://www.python.org/downloads/windows/}
  \bibitem{g07r06} Astral, \emph{python-build-standalone}: relocatable CPython builds.
        \url{https://github.com/astral-sh/python-build-standalone}
\end{reflist}
\end{frame}

\end{document}
